\documentclass{jsarticle}
\usepackage{amsmath}
\usepackage{amssymb}
\begin{document}
\title{整数論から代数、そして幾何学}
\author{Masaaki Yamaguchi}
\date{}
\maketitle



     ある整数をaとおく。このaを$\log x$で割ると 
  $$ x = {a \over \log x}$$
  で求まり、ゼータ関数のxが導かれる。
  このxが、
  $$ a = x \log x $$
  とシャノンのエントロピー値として、
  $$ \int a dx = \int x \log x dx $$
  と、コルモゴロフ・シナイの値にもなる。
  これが、別の道でもある、求まるゼータ関数の式にもなる。
  この式から、
  $$　C = \int({{\int{1 \over x^s}dx} + \log x})\mathrm{dvol} $$
  このオイラーの定数を多様体積分すると
  $$ \int C dx_m = {d \over df}{\int \int{1 \over (y \log y)^{1 \over 2}}}dy_m
  + \int \log x \mathrm{dvol} $$
  となり、hartshorn conjectureの多世界理論となり、いろんな式に分岐していく。


    始めのxを見つけにくいので、オイラーの定数の多様体積分の分岐の式達がある。
    要するに、xの式のヒントが幾多もある。

    ガンマ関数を大域的微分多様体として共変微分すると、
    オイラーの定数を多様体積分した式と合流する。
    次元を多様体を基底群として、オイラーの定数とヒッグス場の式を加群すると、
    一般座標変換を大域的微分方程式として求めた結果と同型ともなる。
    $$ \int C dx_m + {d \over df}F(x) = \Gamma^{\gamma^{'}} $$
    $$ {d \over d\gamma}\Gamma = \Gamma^{\gamma^{'}} $$
    これらの式がJones多項式となり、経済理論をReco level theoryとしての
    一般周期関数ともなる。
    $$ 2i\sin (i x \log x) = e^{f} + e^{-f} $$

    これらから、素数分布論が$ x = {a \over \log x} $が基底群を成して、
    数として、オイラーの定数から、幾多の式になり、Jones多項式を最大数
    と最少数の間を一般周期関数として数学を形成して、整数論が群論となり、
    代数、解析、幾何学を成り立たせている。

















  


\end{document}

